% =============================================================================
%  COMPONENT COOKBOOK  --  not part of your report
%  ---------------------------------------------------------------------------
%  Copy any block below into sections/ and adapt it.
%  To preview: set this file as the main document in Overleaf, compile, then
%  switch the main document back to main.tex.
% =============================================================================
\documentclass[10pt, a4paper]{article}

\usepackage{preamble}
% =============================================================================
%  METADATA  --  the only content you edit for a new report
% =============================================================================

% --- Organisation -----------------------------------------------------------
% Replace these three lines and images/logo.png to reuse the template for any
% other university, faculty, company or club. Leave a line empty to hide it.
\newcommand{\orglineA}{VIETNAM NATIONAL UNIVERSITY HO CHI MINH CITY}
\newcommand{\orglineB}{HO CHI MINH CITY UNIVERSITY OF TECHNOLOGY}
\newcommand{\orglineC}{FACULTY OF COMPUTER SCIENCE AND ENGINEERING}

% --- Report -----------------------------------------------------------------
\newcommand{\doctype}{PROJECT REPORT}          % PROJECT REPORT | ASSIGNMENT | LAB REPORT
\newcommand{\coursename}{DATABASE SYSTEMS}
\newcommand{\coursecode}{CO2014}
\newcommand{\reporttitle}{ASSIGNMENT 1}
\newcommand{\instructor}{Vu Ngoc Tu, CSE-HCMUT}
\newcommand{\locationdate}{Ho Chi Minh City, March, 2026}

% --- Running header (keep it short, one line) -------------------------------
\newcommand{\reporthead}{CO2014 --- Database Systems}


\begin{document}

\tableofcontents
\clearpage

% =============================================================================
\section{Figures}
% =============================================================================

\begin{figure}[H]
    \centering
    \begin{tikzpicture}
        \node[draw=BrandPrimary, fill=BrandLight, rounded corners,
              minimum width=8cm, minimum height=4cm, align=center,
              text=BrandPrimary] {Single figure placeholder};
    \end{tikzpicture}
    \caption{Describe what the reader should notice, not that it is a figure.}
    \label{fig:cookbook-single}
\end{figure}

\begin{figure}[H]
    \centering
    \begin{subfigure}{0.45\linewidth}
        \centering
        \begin{tikzpicture}
            \node[draw=BrandAccent, fill=BrandAccent!8, rounded corners,
                  minimum width=4.5cm, minimum height=3cm, align=center,
                  text=BrandAccent] {(a)};
        \end{tikzpicture}
        \caption{Left panel}
        \label{fig:cookbook-left}
    \end{subfigure}
    \hfill
    \begin{subfigure}{0.45\linewidth}
        \centering
        \begin{tikzpicture}
            \node[draw=BrandPrimary, fill=BrandLight, rounded corners,
                  minimum width=4.5cm, minimum height=3cm, align=center,
                  text=BrandPrimary] {(b)};
        \end{tikzpicture}
        \caption{Right panel}
        \label{fig:cookbook-right}
    \end{subfigure}
    \caption{Side-by-side panels.}
    \label{fig:cookbook-pair}
\end{figure}

Refer to them as Figure~\ref{fig:cookbook-pair} or \autoref{fig:cookbook-single}.

% =============================================================================
\section{Tables}
% =============================================================================

\begin{table}[H]
    \centering
    \caption{A booktabs table: no vertical rules, clear header.}
    \label{tab:cookbook-booktabs}
    \begin{tabular}{lrr}
        \toprule
        \textbf{Item} & \textbf{2025} & \textbf{2026} \\
        \midrule
        Revenue       & 1\,200        & 1\,480        \\
        Cost          & 740           & 810           \\
        Margin        & 460           & 670           \\
        \bottomrule
    \end{tabular}
\end{table}

\begin{table}[H]
    \centering
    \caption{A full-width table whose text columns wrap automatically.}
    \label{tab:cookbook-wide}
    \begin{tabularx}{\textwidth}{lX r}
        \toprule
        \textbf{Stage} & \textbf{Description} & \textbf{Weeks} \\
        \midrule
        Analysis       & Collect requirements, confirm scope and constraints with stakeholders & 2 \\
        Design         & Model the data, define interfaces and review alternatives & 3 \\
        Build          & Implement, test and document the result & 5 \\
        \bottomrule
    \end{tabularx}
\end{table}

% =============================================================================
\section{Lists}
% =============================================================================

\begin{itemize}
    \item First-level item
    \begin{itemize}
        \item Second-level item
        \item Another one with \textbf{emphasis}
    \end{itemize}
    \item Back to the first level
\end{itemize}

\begin{enumerate}
    \item Numbered step
    \item Next step
\end{enumerate}

\begin{description}
    \item[Term] Definition of the term.
    \item[Another term] Its definition.
\end{description}

% =============================================================================
\section{Theorem-like environments}
% =============================================================================

\begin{definition}[Access control]
    A mechanism that restricts which subjects may perform which operations on
    which objects.
\end{definition}

\begin{theorem}
    If every transaction is serialisable, then every schedule is equivalent to
    a serial schedule.
\end{theorem}

\begin{problem}
    Given the schema in Table~\ref{tab:cookbook-booktabs}, normalise the
    relation to third normal form.
\end{problem}

\begin{solution}
    State the steps, then the result.
\end{solution}

% =============================================================================
\section{Code}
% =============================================================================

\begin{lstlisting}[language=Python, caption={A numbered, syntax-highlighted listing.}]
def moving_average(values, window):
    """Return the simple moving average."""
    out = []
    for i in range(len(values) - window + 1):
        chunk = values[i:i + window]
        out.append(sum(chunk) / window)
    return out
\end{lstlisting}

Inline code: \codeinline{pip install -r requirements.txt}.

% =============================================================================
\section{Citations and cross-references}
% =============================================================================

Narrative citation: \textcite{sample-article} argues that examples should be
reproducible. Parenthetical citation: \parencite{sample-book}. Online source
with an access date: \parencite{sample-online}. See
Section~\ref{sec:cookbook-equations} for the equations.

% =============================================================================
\section{Equations}
\label{sec:cookbook-equations}
% =============================================================================

\begin{equation}
    \bar{x} = \frac{1}{n}\sum_{i=1}^{n} x_i
    \label{eq:cookbook-mean}
\end{equation}

Equation~\eqref{eq:cookbook-mean} is referenced by number only.

% =============================================================================
\section{Callout box}
% =============================================================================

\begin{tcolorbox}[colback=BrandLight, colframe=BrandPrimary,
                  boxrule=0.6pt, arc=3pt, title=Note]
    Use a callout for a single takeaway, a warning, or a definition that should
    not be missed while reading.
\end{tcolorbox}

% =============================================================================
\section{Landscape page}
% =============================================================================

\begin{landscape}
\subsection{A wide table on a rotated page}

\begin{table}[H]
    \centering
    \caption{Very wide comparison that does not fit in portrait.}
    \label{tab:cookbook-landscape}
    \begin{tabularx}{\linewidth}{lXXXX}
        \toprule
        \textbf{Criterion} & \textbf{Option A} & \textbf{Option B} & \textbf{Option C} & \textbf{Option D} \\
        \midrule
        Cost               & Low               & Medium            & High              & Medium            \\
        Effort             & Low               & High              & Medium            & Low               \\
        Risk               & Medium            & Low               & Low               & High              \\
        \bottomrule
    \end{tabularx}
\end{table}
\end{landscape}

% =============================================================================
\section{Appendix entries}
% =============================================================================

After \verb|\appendix| in main.tex, every \verb|\section| becomes A, B, C and
can be referenced with \verb|\ref{app:supplementary}|. See
\verb|appendices/A-appendix.tex| for a ready-made example.

\end{document}
